\documentclass[12pt,a4paper]{article}
\usepackage[utf8]{inputenc}
\usepackage[margin=2.5cm]{geometry}

\title{PE5 -- Vista previa individual\\Resumen (español) y Abstract (inglés)}
\author{Nieves Sánchez, Jimmy Samuel}
\date{}

\begin{document}
\maketitle



\section*{Resumen}

MundiPets es una plataforma digital para la adopción y cruza responsable de
mascotas, que centraliza la información de los animales y fortalece la
confiabilidad de los datos publicados frente a procesos informales, opacos y
propensos al fraude, documentados en PE1-PE2. El equipo aplicó un proceso de
ingeniería de requisitos conforme a ISO/IEC/IEEE 29148:2018, con inspección
formal Fagan y gestión de cambios mediante RFC, consolidando un ERS versión
4.0 con 61 requisitos (27 RF, 16 RNF, 9 RD, 9 RL), 13 casos de uso, una
matriz de trazabilidad sin huérfanos y tres componentes de inteligencia
artificial con requisitos de equidad medibles. La auditoría de calidad
mostró resultados dispares entre las seis métricas: la corrección tras la
inspección alcanzó el valor óptimo (M6 = 0,00), pero la completitud
estructural de atributos por tipo de requisito quedó muy por debajo del
umbral exigido (M1a = 44,26\,\%), por heterogeneidad de plantilla entre
categorías, no por errores de contenido. Se concluye que la verificación
formal, aunque eficaz para el contenido que audita, no cubre por diseño la
consistencia estructural entre secciones ni entre los requisitos y los
modelos UML, por lo que el equipo adopta como práctica extender la
verificación cruzada más allá del alcance de la inspección original.

\textbf{Palabras clave:} ingeniería de requisitos, ISO/IEC/IEEE 29148:2018,
trazabilidad end-to-end, auditoría de calidad, requisitos de inteligencia
artificial, MundiPets.

\bigskip

\section*{Abstract}

MundiPets is a digital platform for responsible pet adoption and breeding
that centralizes animal information and strengthens the reliability of
published data, addressing informal, opaque, and fraud-prone processes
documented through interviews and questionnaires with owners, adopters, and
shelters during PE1-PE2. The team applied a requirements engineering
process aligned with ISO/IEC/IEEE 29148:2018, including formal Fagan
inspection and change management through RFCs and a Change Control Board,
resulting in a final SRS version 4.0 with 61 requirements (27 functional,
16 non-functional, 9 design constraints, 9 legal), 13 use cases, an
end-to-end traceability matrix with no orphan requirements, and three
artificial-intelligence components with measurable fairness requirements.
The quality audit showed uneven results across the six metrics:
post-inspection correctness reached the optimal value (M6 = 0.00), while
structural attribute completeness by requirement type fell far below the
required threshold (M1a = 44.26\,\%), due to template heterogeneity across
categories rather than content errors. The team concludes that formal
verification, while effective for the content it audits, is not designed by
scope to cover structural consistency across sections or between
requirements and UML models, leading the team to adopt cross-artifact
verification beyond the original inspection scope as standard practice
going forward.

\textbf{Keywords:} requirements engineering, ISO/IEC/IEEE 29148:2018,
end-to-end traceability, quality audit, artificial intelligence
requirements, MundiPets.

\end{document}
